\section{Discussion}
\label{sec:discussion}

The central finding of this work is that existing kinase inhibitor selectivity
definitions do not measure the same quantity.
The ratio definition and the distribution-based family (S-score, entropy, Gini)
produce substantially different compound rankings (cross-family
$r = 0.14$--$0.62$ across four datasets, versus within-family $r > 0.72$), and
the sources of disagreement are mechanistically interpretable rather than random.
This has practical consequences: a compound that ratio-based analysis designates
as selective may be designated as moderately promiscuous by entropy-based
analysis, and vice versa, depending on the shape of its binding profile.
Go/no-go decisions in early drug discovery are routinely made on the basis of
selectivity assessments without a principled choice of selectivity definition.

The clustering of definitions into two families
reflects a difference in what is being measured.
Distribution-based definitions ask: how concentrated is the compound's binding
activity across the full profiled kinome?
The ratio definition asks: how much more potently does the compound bind its
primary target than a specified off-target?
The selectivity measure discrepancy is most pronounced 
for compounds with one clearly dominant target that also
bind many kinases at moderate affinity (Type~3 profiles in our taxonomy).
For such compounds, ratio-based methods report high selectivity while
distribution-based methods report moderate promiscuity.
The required properties proposed in the Required Properties section are an
attempt to express what it means for a compound to be selective in a way that is
precise enough to assess candidate selectivity definitions formally.

This coexistence has a historical explanation.
In classical pharmacology and analytical chemistry, selectivity is inherently
pairwise. IUPAC formalizes it as a selectivity coefficient, the ratio of
binding constants for an analyte and an interferent\cite{IUPAC2001}, and
ratio measures entered kinase profiling on that basis.
Kinome-wide profiling, however, poses a distributional question: how
concentrated is a compound's binding across hundreds of kinases?
Entropy and Gini were introduced because ratio measures are ill-suited to it, since
they discard all but two points of a 300--500 kinase profile. 


\subsection{Practical Recommendations}

Pending a formal resolution, three practical recommendations follow from our
findings.
First, D1 (reliability threshold) should be adopted as a reporting
standard.
Compounds with no detected binding above the assay detection floor should not
receive selectivity scores, or should have their scores explicitly flagged as
unreliable.
Of the 222 Klaeger compounds, 16 (7.2\%) fall into this category and currently
receive scores that are indistinguishable from noise.

Second, selectivity assessments should report the sensitivity of rankings to
parameter choice rather than a single score at a single parameter value.
Our analysis shows that rank standard deviation across the parameter sweep is a
compound-specific, interpretable quantity that identifies which compounds have
robust selectivity assessments and which do not.
This is analogous to reporting confidence intervals in statistical estimation.

Third, when comparing selectivity across studies that use different definitions,
the comparison should be restricted to the subset of compounds and parameter
values for which the definitions agree.
The high within-family correlations ($r > 0.72$) suggest that S-score, entropy,
and Gini are interchangeable for most purposes when parameters are chosen
consistently. This is an empirical observation rather than a preset adequacy
bar: $r = 0.72$ is the lowest within-family correlation we observed (Klaeger
S-score vs.\ Gini, Table~\ref{tab:correlations}), and it measures a different
property than the $\rho > 0.90$ criterion used in the panel-size analysis
below. The former measures how similarly two distinct definitions rank compounds
on the full panel; the latter measures how similarly the same definition ranks
compounds on a subpanel versus the full panel. The
within-family agreement is measured on the complete kinome, and the panel-size
analysis shows that the Gini coefficient can diverge from the entropy and
S-score on the smaller panels used in practice (on Klaeger it requires roughly
290 of 343 kinases to recover its full-panel ranking, against 110 for
entropy and the S-score). The three measures should therefore be treated as
interchangeable only when compared on panels of similar, near-complete coverage.
Ratio-based measures, in turn, should not be compared directly with
distribution-based measures at any panel size without acknowledging the
conceptual distinction.

\subsection{Limitations}

The analysis covers four profiling datasets spanning three assay technologies
and two readout types, but these still represent a fraction of the chemical and
kinase space relevant to drug discovery.
The assay technologies are not directly comparable in absolute terms --- the
active-compound fractions differ across datasets. As a consequence,
our analysis
relies on within-dataset rank comparisons rather than cross-dataset absolute
scores.
The consistency of the two-family structure and the instability predictors
across all four datasets indicates that the findings are not artifacts of a
single assay platform. However, findings should still be interpreted in light of
these technological differences.
Our comparison is also restricted to four widely used, historically
established definitions. Bosc et al.'s window score and ranking
score\cite{Bosc2017} were not included, but a check on the same four datasets
confirms both correlate with the ratio family (Spearman $r = 0.80$--$0.96$)
rather than forming a distinct third cluster.
The CATDS score\cite{Klaeger2017}, historically classified alongside the
ratio-based measures, was also not implemented directly. Unlike $R(k)$, CATDS
normalizes each target's binding reduction by the sum of reductions across the
entire panel rather than comparing to a single off-target, a construction it
shares with the distribution-based family rather than with $R(k)$. A proxy
computed from our baseline-subtracted affinity profiles, in place of the raw
concentration-response data CATDS requires, correlates more strongly with Gini
and entropy than with $R(k)$ on Klaeger, consistent with this structural
reading. We did not add it as a fifth measure: the distribution family is
already represented by three members, and CATDS appears to be an additional
instance of the same underlying construct rather than an independent
perspective (see repository for results).

A further axis of platform dependence that this analysis does not examine is
cell-based versus cell-free profiling. Recent work shows that selectivity
classifications for the same compounds can differ substantially between
biochemical binding panels and cellular target-engagement assays. For example, several
type II (DFG-out-binding) inhibitors show cellular interactions that
cell-free assays miss\cite{Binder2026}. A plausible contributor, not established
by that work, is that cell-free panels often use truncated or constitutively
active kinase-domain constructs, which cannot adopt the DFG-out state.

We also examined whether in vitro selectivity predicts clinical toxicity,
correlating the selectivity rankings of the FDA-approved Klaeger compounds with
FAERS adverse-event rates ($n = 46$) and FDA-label discontinuation rates
($n = 33$).
None of the rate-based outcomes were significantly associated with selectivity
under any definition (all $|r| < 0.17$, all $p > 0.27$).
The only nominal signal is a weak correlation between raw FAERS report counts and
the ratio ranking ($r = 0.31$, $p = 0.035$), which most likely reflects
prescription volume rather than toxicity.
We do not read the absence of a rate-based association as evidence that
selectivity and toxicity are unrelated. FAERS report counts mostly track how
often a drug is prescribed, and label discontinuation depends on the patient
population treated for each indication, so neither can isolate the effect of
selectivity.
A proper test would require indication-stratified, treatment-related
adverse-event rates from controlled trials with matched comparators, which are
not available in standardized form across these drugs.

The four properties proposed above are necessary
conditions motivated by empirical failure modes.
The candidate measure shows that they are jointly satisfiable, but they are not
proven to be sufficient to uniquely characterize a selectivity measure, and we do
not show that they single out a measure of the candidate's form.
The satisfaction assessments in Table~\ref{tab:desiderata} are based on
analysis of the functional forms of each definition together with the empirical
checks of this study. A formal proof of satisfaction or violation for each
property--definition pair is left for future work.
